% ============================================================
%  DevFest GDGoC Japan 2026 Hackathon - Submission Document
%  Compile with: pdflatex contour_submission.tex  (run twice)
%  Needs the icons/ folder next to this file (architecture diagram logos).
%  Red [TBC: ...] marks = not confirmed from the repo; fill in before submitting.
% ============================================================
\documentclass[11pt,a4paper]{article}

\usepackage[T1]{fontenc}
\usepackage[utf8]{inputenc}
\usepackage{lmodern}
\usepackage[margin=25mm]{geometry}
\usepackage{amssymb}
\usepackage{array}
\usepackage{tabularx}
\usepackage{booktabs}
\usepackage{enumitem}
\usepackage{xcolor}
\usepackage{titlesec}
\usepackage{graphicx}
\usepackage{tikz}
\usepackage{float}
\usepackage[hidelinks]{hyperref}

\usetikzlibrary{positioning,arrows.meta,fit,backgrounds,calc,svg.path}

% ---------- colours ----------
\definecolor{gblue}{HTML}{1A73E8}
\definecolor{gframe}{HTML}{BDC1C6}
\definecolor{ghead}{HTML}{F1F3F4}
\definecolor{gtext}{HTML}{202124}

% ---------- section styling ----------
\titleformat{\section}{\large\bfseries\color{gtext}}{\thesection.}{0.6em}{}
\titleformat{\subsection}{\normalsize\bfseries}{\thesubsection.}{0.6em}{}
\titlespacing*{\section}{0pt}{16pt}{8pt}
\titlespacing*{\subsection}{0pt}{10pt}{5pt}

% ---------- helpers ----------
\newcommand{\checkedbox}{\makebox[1.2em]{$\boxtimes$}}
\newcommand{\emptybox}{\makebox[1.2em]{$\square$}}
% not confirmed yet: shows in red so it is easy to find before submitting
\newcommand{\tbc}[1]{\textcolor{red}{[TBC: #1]}}

% a light framed block, used for the free-text answer boxes
\newsavebox{\answerboxcontent}
\newenvironment{answerbox}
  {\par\vspace{4pt}\setlength{\fboxsep}{10pt}\setlength{\fboxrule}{0.4pt}%
   \begin{lrbox}{\answerboxcontent}%
   \begin{minipage}{\dimexpr\textwidth-2\fboxsep-2\fboxrule\relax}}
  {\end{minipage}\end{lrbox}%
   \noindent\fbox{\usebox{\answerboxcontent}}\par\vspace{4pt}}

% table column types
\newcolumntype{L}[1]{>{\raggedright\arraybackslash}p{#1}}
\newcolumntype{Y}{>{\raggedright\arraybackslash}X}

\renewcommand{\arraystretch}{1.35}

% bracketed part headings, e.g. [Required Items]
\newcommand{\partheading}[1]{%
  \vspace{14pt}\noindent{\large\bfseries [#1]}\par\vspace{4pt}%
  \noindent\rule{\textwidth}{0.8pt}\par\vspace{2pt}}

% ---------- architecture diagram (same as architecture.tex) ----------
\graphicspath{{icons/}}

%------------------------------- Colours --------------------------------------
\definecolor{stepC}{HTML}{1A3D8F}      % dark blue step numbers
\definecolor{ink2}{HTML}{4A524C}
% box fills / borders by group
\definecolor{phoneF}{HTML}{E6F4EA}  \definecolor{phoneD}{HTML}{34A853}
\definecolor{fbF}{HTML}{FFF3E0}     \definecolor{fbD}{HTML}{FFA000}
\definecolor{gcpF}{HTML}{E8F0FE}    \definecolor{gcpD}{HTML}{4285F4}
\definecolor{aiF}{HTML}{F3E8FD}     \definecolor{aiD}{HTML}{8E75B2}
\definecolor{extF}{HTML}{F1F3F4}    \definecolor{extD}{HTML}{9AA0A6}

%------------------------------- Logos ----------------------------------------
% Firebase and Android: PNG files in icons/. Gemini: vector path (simple-icons, CC0).
\newcommand{\svglogo}[3][0.028]{%
  \tikz[x=#1cm, y=-#1cm, baseline={([yshift=-3pt]current bounding box.center)}]
    \fill[#2, even odd rule] svg {#3};}
\newcommand{\firebaselogo}{\raisebox{-0.5\height}{\includegraphics[width=0.68cm]{firebase}}}
\newcommand{\androidlogo}{\raisebox{-0.5\height}{\includegraphics[height=0.68cm]{android}}}
% small inline platform marks for the User title
\newcommand{\platforms}{\hspace{3pt}\raisebox{-1pt}{\includegraphics[height=0.3cm]{android}}%
  \hspace{2pt}\raisebox{-1pt}{\includegraphics[height=0.3cm]{apple}}}
\newcommand{\mlkitlogo}{\raisebox{-0.5\height}{\includegraphics[width=0.6cm]{mlkit}}}
% Scheduler + Workflows: two icons stacked
\newcommand{\jobslogos}{\begin{tabular}{@{}c@{}}\includegraphics[width=0.45cm]{cloud_scheduler}\\[1pt]
  \includegraphics[width=0.45cm]{workflows}\end{tabular}}
\newcommand{\geminilogo}{\svglogo{aiD}{M11.04 19.32Q12 21.51 12 24q0-2.49.93-4.68.96-2.19 2.58-3.81t3.81-2.55Q21.51 12 24 12q-2.49 0-4.68-.93a12.3 12.3 0 0 1-3.81-2.58 12.3 12.3 0 0 1-2.58-3.81Q12 2.49 12 0q0 2.49-.96 4.68-.93 2.19-2.55 3.81a12.3 12.3 0 0 1-3.81 2.58Q2.49 12 0 12q2.49 0 4.68.96 2.19.93 3.81 2.55t2.55 3.81}}
% Simple drawn icons for the non-product boxes
\newcommand{\usericon}{\tikz[baseline={([yshift=-3pt]current bounding box.center)}]{%
  \fill[phoneD] (0,0.47) circle (0.14);
  \fill[phoneD, rounded corners=3pt] (-0.25,0) rectangle (0.25,0.3);}}
\newcommand{\docicon}{\tikz[baseline={([yshift=-3pt]current bounding box.center)}]{%
  \draw[extD, line width=0.8pt, fill=white] (0,0) -- (0.5,0) -- (0.5,0.5) -- (0.36,0.64) -- (0,0.64) -- cycle;
  \foreach \y in {0.15,0.27,0.39} \draw[extD, line width=0.6pt] (0.1,\y) -- (0.4,\y);}}
\newcommand{\png}[1]{\raisebox{-0.5\height}{\includegraphics[width=0.68cm]{#1}}}

% Box content: icon on the left, title + subtitle on the right
\newcommand{\svc}[3]{%
  \begin{tabular}{@{}m{0.6cm}@{\hspace{5pt}}>{\raggedright\arraybackslash}m{2.85cm}@{}}
    \centering #1 & {\footnotesize\bfseries #2}\par\vspace{1pt}{\scriptsize\color{ink2}#3}
  \end{tabular}}

% Step label: dark-blue bold number, then the text
\newcommand{\step}[2]{{\color{stepC}\bfseries\footnotesize #1}~%
  \begin{tabular}[c]{@{}l@{}}#2\end{tabular}}


% ============================================================
\begin{document}

% ------------------------------------------------------------
% Title block
% ------------------------------------------------------------
\begin{center}
  {\LARGE\bfseries DevFest GDGoC Japan 2026 Hackathon}\\[6pt]
  {\large Submission Document}
\end{center}
\vspace{6pt}

\noindent\textbf{Team Information}\par\vspace{4pt}
\noindent
\begin{tabularx}{\textwidth}{@{}L{4.2cm}Y@{}}
\toprule
Venue                  & \tbc{venue} \\
Team Name              & \tbc{team name} \\
Member                 & \tbc{member names} \\
Application Name       & ``Contour'' \\
Category               & \tbc{category, e.g.\ disaster prevention} \\
GitHub Repository URL  & \url{https://github.com/dpbelisario/Devfest2026} \\
\bottomrule
\end{tabularx}

\partheading{Required Items}

% ------------------------------------------------------------
\section{Project Overview}

``Contour'' is a phone app that helps people in Nagoya find the nearest official
evacuation site and get there, in their own language. It shows the city's designated
emergency evacuation sites on a map, points the way with 3D AR arrows on the camera view,
and answers questions through a Gemini assistant.

The same app also works as a campus guide for Nagoya University students (Student mode),
including directions to a classroom inside a building.

\begin{itemize}[leftmargin=1.4em,itemsep=2pt,topsep=4pt]
  \item \textbf{Supported languages:} English, Japanese, Chinese and Korean
  \item \textbf{Platforms:} Android app, iOS app shell and any phone browser
        (all show the same web page)
  \item \textbf{Data:} 1,735 official evacuation sites for Nagoya City (GSI open data)
        and Japan Meteorological Agency (JMA) warnings, refreshed automatically
  \item \textbf{Reliability:} Only official designations are shown. Opening status and
        capacity stay ``unknown'' instead of being guessed, and the assistant points
        users to official instructions.
\end{itemize}

% ------------------------------------------------------------
\section{Problems Addressed}

\subsection{Target Users}
\begin{itemize}[leftmargin=1.4em,itemsep=2pt,topsep=4pt]
  \item Foreign residents, visitors and international students in Nagoya
  \item Nagoya University students looking for a building or classroom
  \item Anyone who needs to find the nearest evacuation site quickly
\end{itemize}

\subsection{Problems and Rationale}
Official evacuation-site information exists, but it is published as Japanese map
services and data files, and it does not say how to get to a site from where you stand.
Emergency alerts on Japanese phones also arrive in Japanese.

For people who do not read Japanese well, this means losing time at exactly the moment
it matters most. Contour puts the nearest official sites, the direction to walk and what
to do now on one screen, in the user's language.

\subsection{Differences from Existing Approaches}

\noindent
\begin{tabularx}{\textwidth}{@{}L{3.0cm}YY@{}}
\toprule
\textbf{Aspect} & \textbf{Official hazard maps and general map apps} & \textbf{Contour} \\
\midrule
Finding a site      & Search a map service or data file, mostly in Japanese.            & Nearest official sites on a map, filtered by disaster type, with distance and direction. \\
Getting there       & Read a 2D map.                                                    & 3D AR arrows on the camera view, a walking-route preview, and indoor directions to the nearest exit (demo). \\
Emergency alerts    & The phone's alert text, in Japanese.                              & On Android the app opens with the alert, what to do now and the nearest suitable site, in the app language. \\
Questions           & None.                                                             & Gemini assistant that answers from the real site and warning data. \\
Languages           & Mostly Japanese.                                                  & English, Japanese, Chinese and Korean, with on-device translation. \\
\bottomrule
\end{tabularx}

% ------------------------------------------------------------
\section{Use of Google Technologies}

\subsection{Google Technologies Used (select at least one)}

\noindent
\begin{tabularx}{\textwidth}{@{}YY@{}}
\emptybox~App Engine                  & \checkedbox~Cloud Run \\
\emptybox~Compute Engine              & \emptybox~Firebase App Hosting \\
\emptybox~Google Kubernetes Engine (GKE) & \checkedbox~Firebase Hosting \\
\end{tabularx}

\subsection{Examples of Google Technology Use}

\noindent
\begin{tabularx}{\textwidth}{@{}L{2.9cm}YYY@{}}
\toprule
\textbf{Technology} & \textbf{Use} & \textbf{Reasons / Benefits} & \textbf{Drawbacks} \\
\midrule
Cloud Run & Hosts the backend API and the web page; a second, private service runs the data-import jobs. & One container image for both services; scales down to zero when idle. & Cold starts can delay the first request. Container disk is not permanent storage. \\
Firebase Hosting & Public web address that forwards every request to Cloud Run. & One HTTPS address for the Android app, iOS app and browsers. & Adds a forwarding step that has to be configured. \\
Gemini (Vertex AI) & ``Ask Contour'' assistant with function calling: finds sites, reads warnings, shows sites on the map or starts AR. & Natural replies in four languages, grounded in our own data. Usually well under a cent per question. & Needs rate limits (8 per minute per user) and strict instructions so it does not invent sites or warnings. \\
Cloud Firestore & Stores the imported evacuation sites, JMA bulletins and a warning outbox. & Serverless; the API keeps a short in-memory copy, so most requests cost no reads. & Charged per read and write, which is why caching was needed. \\
Cloud Scheduler + Cloud Workflows & Run the imports: sites every hour, JMA warnings every minute. & Retries and a visual step graph without running our own server. & Workflow definitions can drift from the repository if edited in the console. \\
ML Kit Translation + Android speech & Offline translation and voice input on Android. & Works without the internet, on the phone. & Each language model must be downloaded first; Android only. \\
Cloud Build + Firebase App Distribution & Build the Android app on each push and send it to testers. & Teammates update the app with one tap. & iOS builds still need Apple signing. \\
\bottomrule
\end{tabularx}

\subsection{Why We Chose Google Technologies}

\begin{answerbox}
We needed a working, multilingual app within a short hackathon. Google Cloud and Firebase
let us put the website, the API, the database, the scheduled data imports and the AI
assistant in one project, without managing servers.

Gemini's function calling was the key piece. Instead of letting the model answer from
memory, the assistant calls our own tools to read the official site list and the latest
JMA warnings, so its answers stay tied to real data.
\end{answerbox}

% ------------------------------------------------------------
\section{Architecture Diagram}

\begin{figure}[H]
\centering\sffamily
\hyphenpenalty=10000
\exhyphenpenalty=10000
\sloppy

\begin{tikzpicture}[
  box/.style={rounded corners=4pt, line width=0.8pt, inner sep=5pt,
              minimum width=3.6cm, minimum height=1.3cm},
  phone/.style={box, fill=phoneF, draw=phoneD},
  fb/.style={box, fill=fbF, draw=fbD},
  gcp/.style={box, fill=gcpF, draw=gcpD},
  core/.style={box, fill=gcpF, draw=gcpD, line width=1.6pt},
  ai/.style={box, fill=aiF, draw=aiD},
  ext/.style={box, fill=extF, draw=extD},
  flow/.style={-{Stealth[length=2.4mm]}, draw=stepC!60, line width=0.9pt},
  lbl/.style={font=\scriptsize, text=ink2, inner sep=2pt, align=left}
]

% ---------------- Boxes (top to bottom) ----------------
\node[phone] (user)   at (0,0)
  {\svc{\usericon}{User\platforms}{Android / iOS app showing one shared web page with map, 3D AR arrows, indoor nav.}};
\node[phone] (device) at (-5.5,0)
  {\svc{\mlkitlogo}{ML Kit + Speech}{Offline translation and voice input, AR uses camera + compass}};

\node[fb]    (host)   at (0,-4.5)
  {\svc{\firebaselogo}{Firebase Hosting}{Public web address, forwards every request to Cloud Run}};
\node[core]  (run)    at (0,-8.8)
  {\svc{\png{cloud_run}}{Cloud Run}{Serves the web page and API for nearest sites, alerts, assistant}};
\node[ai]    (gemini) at (-5.5,-8.8)
  {\svc{\geminilogo}{Gemini}{Vertex AI. Replies in English, Japanese, Chinese or Korean with map or AR actions}};
\node[gcp]   (store)  at (0,-13.1)
  {\svc{\png{firestore}}{Cloud Firestore}{Evacuation sites, JMA alerts, warning outbox}};
\node[gcp]   (jobs)   at (5.5,-13.1)
  {\svc{\jobslogos}{Scheduler + Workflows}{Runs import jobs and checks data before saving}};

\node[ext]   (gov)    at (5.5,-17.0)
  {\svc{\docicon}{Official Open Data}{Nagoya evacuation-site list, JMA warning feeds}};

% ---------------- Group frames ----------------
\begin{scope}[on background layer]
  \node[draw=phoneD!60, dashed, rounded corners=6pt, inner sep=8pt, fit=(user)(device)] (phoneG) {};
  \node[anchor=south west, font=\scriptsize\bfseries, text=phoneD!80!black, inner sep=2pt]
        at (phoneG.north west) {On the phone};
  \node[draw=gcpD!50, dashed, rounded corners=6pt, inner sep=8pt, fill=gcpF!35,
        fit=(host)(gemini)(run)(store)(jobs)] (cloud) {};
  \node[anchor=north east, font=\scriptsize\bfseries, text=gcpD!80!black, inner sep=4pt]
        at (cloud.north east) {Google Cloud / Firebase};
\end{scope}

% ---------------- User request: 1 -> 4 (straight down the middle) ----------------
\draw[flow] (user.south)  -- node[lbl, right]
  {\step{1}{Open the map or ask a question\\(location + language sent)}} (host.north);
\draw[flow] (host.south)  -- node[lbl, right]
  {\step{2}{HTTPS request for sites, alerts\\or assistant chat}} (run.north);
\draw[flow] (run.south)   -- node[lbl, right]
  {\step{3}{Read sites + alerts\\(cached 5 min / 30 s)}} (store.north);
\draw[flow] (run.west)    -- node[lbl, above]
  {\step{4}{Ask\\Gemini}} (gemini.east);

% On the phone, no server involved
\draw[flow, densely dotted, draw=phoneD, {Stealth[length=2mm]}-{Stealth[length=2mm]}]
  (device.east) -- node[lbl, above] {on phone} (user.west);

% ---------------- Background import: A -> B ----------------
\draw[flow, dashed, draw=extD] (gov.north) -- node[lbl, left, pos=0.45]
  {\step{A}{Fetch sites every hour,\\warnings every minute}} (jobs.south);
\draw[flow] (jobs.west) -- node[lbl, above] {\step{B}{Save}} (store.east);

\end{tikzpicture}

\caption{System architecture of Contour. Steps 1--4 follow one user request from top to bottom,
results return the same way. Steps A--B run on a schedule (evacuation sites hourly, JMA warnings every minute).}
\label{fig:contour-architecture}
\end{figure}

% ------------------------------------------------------------
\section{Deployed Project URL}

\begin{answerbox}
Application URL: \url{https://gen-lang-client-0467526025.web.app/app/}\\[4pt]
Opens in any phone or desktop browser. The Android app shows the same page; testers
install it through Firebase App Distribution.
\end{answerbox}

\partheading{Optional Items}

% ------------------------------------------------------------
\section{How to Test}

\subsection{Test Account Information}

\noindent
\begin{tabularx}{\textwidth}{@{}L{4.2cm}Y@{}}
\toprule
\textbf{Item} & \textbf{Details} \\
\midrule
URL           & \url{https://gen-lang-client-0467526025.web.app/app/} \\
Account       & Not needed. The app has no login. \\
Android app   & \tbc{App Distribution invite link, or say ``browser only''} \\
\bottomrule
\end{tabularx}

\subsection{Test Procedure}
\begin{enumerate}[leftmargin=1.6em,itemsep=2pt,topsep=4pt]
  \item Open the URL above, ideally on a phone.
  \item Choose a language from the menu in the header (English, Japanese, Chinese or Korean).
  \item In \textbf{Evacuation} mode, check that official Nagoya evacuation sites appear
        on the map. Without location access, a demo position near Nagoya University is used.
  \item Tap a site card to see its details, then tap \textbf{3D AR} and allow the camera
        to see the arrow pointing to the site.
  \item Open \textbf{Ask} and type ``Where is the nearest evacuation site?'' The assistant
        should name nearby official sites and offer to show them on the map.
  \item Switch to \textbf{Student} mode, open \textbf{Ask} and type ``How do I get to
        room 221?'' Step-by-step indoor directions should appear; open \textbf{3D AR} to
        follow them.
  \item (Android app only) Open the menu and tap \textbf{Test alert}. After a few seconds
        the emergency pop-up should open with what to do now and the nearest site.
\end{enumerate}

\subsection{Sample Data}
No account or saved data is needed. The app uses:
\begin{itemize}[leftmargin=1.4em,itemsep=2pt,topsep=4pt]
  \item 1,735 designated emergency evacuation sites for Nagoya City (GSI open data,
        imported automatically every hour)
  \item Live JMA weather, earthquake, tsunami and volcano bulletins (checked every minute)
  \item 16 Nagoya University buildings for Student mode, and a Room 221 indoor demo in
        Engineering Building II
\end{itemize}

% ------------------------------------------------------------
\section{Additional Application and Feature Details}

\subsection{Main Features}

\noindent
\begin{tabularx}{\textwidth}{@{}L{0.8cm}L{3.8cm}Y@{}}
\toprule
\textbf{\#} & \textbf{Feature} & \textbf{Description} \\
\midrule
1 & Nearest evacuation sites & Official Nagoya sites on a map, filtered by disaster type (earthquake, tsunami, flood, landslide, fire and more), with straight-line distance and direction. \\
2 & 3D AR directions & Arrows drawn on the camera view using the phone's compass point toward the chosen site. \\
3 & Ask Contour (Gemini) & Questions in four languages. The assistant reads the real site list and warnings, and can highlight sites on the map or start AR. \\
4 & Emergency pop-up (Android) & When an official emergency alert arrives, Contour opens over the lock screen with what to do now and the nearest suitable site. \\
5 & Indoor navigation (demo) & Guides to the nearest exit (avoiding elevators) or to a classroom, estimating position from steps, compass and air pressure. \\
6 & Student mode & Campus guide for 16 Nagoya University buildings that works offline. \\
7 & Offline translation and voice input & On Android, ML Kit translates site names and text on the phone, and questions can be spoken. \\
\bottomrule
\end{tabularx}

\subsection{Future Plans}
\begin{itemize}[leftmargin=1.4em,itemsep=2pt,topsep=4pt]
  \item Push JMA warnings to phones with Firebase Cloud Messaging (the outbox already exists)
  \item Real floor plans and more accurate indoor positioning for more campus buildings
  \item Verified walking routes that take hazards and closures into account
  \item Distribute the iOS app to testers
\end{itemize}

\partheading{Items Not Scored}

% ------------------------------------------------------------
\section{Where, How and Why We Used AI}

\begin{itemize}[leftmargin=1.4em,itemsep=3pt,topsep=4pt]
  \item \textbf{Development:} \tbc{which AI coding tools the team used and for what.
        The repository includes instructions for AI coding agents (\texttt{AGENTS.md}).}
  \item \textbf{Service features:} Gemini on Vertex AI powers the ``Ask Contour''
        assistant. It calls our own tools to read the official site list and current
        warnings, and replies in the user's language.
\end{itemize}

\tbc{how AI-generated code and answers were reviewed}

% ------------------------------------------------------------
\section{What We Gained and How Much We Grew Through Development}

\tbc{the team's own reflection}

\end{document}
